\documentclass[a4paper,11pt]{article}
\usepackage[utf8]{inputenc}
\usepackage[T1]{fontenc}
\usepackage[english]{babel}
\usepackage[left=2.5cm,right=2.5cm,top=2cm,bottom=2cm]{geometry}
\usepackage{hyperref}
\usepackage{xcolor}
\usepackage{microtype}
\usepackage{lmodern}

\definecolor{primary}{RGB}{0, 51, 102}
\hypersetup{colorlinks=true, urlcolor=primary}

\begin{document}
\noindent
\textbf{Loic Aron MBASSI EWOLO} \\
ENSEA -- Double Diplome Ingenieur \\
\href{mailto:loicaron.mbassiewolo@ensea.fr}{loicaron.mbassiewolo@ensea.fr} \quad (+33) 7 59 52 33 98 \\
\href{https://github.com/Nameless0l}{github.com/Nameless0l} \quad \href{https://mbassiloic.tech}{mbassiloic.tech}

\vspace{1em}
\noindent
Cergy, le 22 mars 2026

\vspace{1em}
\noindent
\textbf{Objet : Candidature -- Data Scientist Corporate AI Squad} \\
\textbf{Reference : Offre d'alternance 12 mois -- Safran}

\vspace{1.5em}
\noindent Madame, Monsieur,

L'aeronautique et la defense constituent des domaines critiques ou les donnees de performance, d'optimisation et de fiabilite determinentt la competitivite et la securite. Ces defis exigent des pipelines de traitement de donnees robustes, des modeles de prediction fiables et une inference temps reel en environnement contraint. Actuellement en stage de recherche a INRIA Saclay (equipe TRIBE, avril a aout 2026), je developpe et deploie des pipelines NLP production pour evaluer la confiance et la transparence dans les applications mobiles, exploitant BERT, les Transformers et le traitement distribue de donnees a grande echelle. Cette experience d'ingenierie donnees appliquee aux defis critiques est directement transferable a vos enjeux d'optimisation et de prediction en aeronautique.

Ma recherche appliquee en deep learning demontre une competence confirmee dans la conception et le deploiement de modeles complexes en production. Au sein de MILA (Institut Quebecois d'IA), j'ai etudie le phenomene de Grokking (generalisation tardive apres suraprentissage) sur des reseaux de neurones, validant des approches theoriques et mathematiques rigoureuses en environnement production. J'ai aussi concepun systeme multi-agents collaboratifs orchestres par Gemini avec RAG pour des recommandations multi-criteres distribuees. Cette competence a architecturer des pipelines complexes et a valider des modeles theoriques en environnement reel correspond a vos besoins en IA pour l'aeronautique.

Lors du UROP ECG (Google DeepMind), j'ai concepun un modele de classification d'anomalies cardiaque a partir d'images d'electrocardiogrammes, travaillant sur le traitement du signal, l'extraction de features et le deep learning medical. J'ai egalement remporte le 1er prix du hackathon IndabaX Africa 2025 en IA Sante avec un projet complet de nettoyage de donnees medicales, deploiement d'API et creation de dashboards. Cette experience du traitement de donnees reelles complexes et du deploiement en environnement production me donne les competences necessaires pour vos projets de data science aeronautique.

Je suis disponible des septembre 2026 pour une alternance de 12 mois et j'apporte une formation de double diplome ENSEA (Systemes Embarques, Signal, IA) associee a une experience de recherche a INRIA. Safran represente un acteur strategique de l'innovation technologique et industrielle francaise, et je suis convaincu que mon experience en IA et traitement de donnees contribuera significativement a vos projets strategiques. Je serais ravi de discuter de ma candidature lors d'un entretien.

\vspace{1.5em}
\noindent Veuillez agreer, Madame, Monsieur, l'expression de mes salutations distinguees.

\vspace{2em}
\noindent \textbf{Loic Aron MBASSI EWOLO}
\end{document}